\section*{Capítulo \ref{ch:b1:p-block-16-18} --- El bloque p II: el oxígeno, los halógenos y los gases nobles}

\begin{solution}{exo:b1:p-block-16-18:1}
\ce{H2S} $-$II; \ce{S8} 0; \ce{SO2} $+$IV; \ce{SO3^2-} $+$IV; \ce{SO4^2-}
$+$VI; \ce{S2O3^2-} $+$II de media; \ce{H2S2O7} $+$VI.
\end{solution}

\begin{solution}{exo:b1:p-block-16-18:2}
Un \omterm{def:b1:p-block-16-18:halogen}{halógeno} oxida los haluros que están debajo de él:
\ce{Cl2 + 2I- -> 2Cl- + I2} y \ce{Cl2 + 2Br- -> 2Cl- + Br2} tienen lugar;
el yodo con el bromuro y el bromo con el cloruro, no.
\end{solution}

\begin{solution}{exo:b1:p-block-16-18:3}
\ce{SO2}: S con un \omterm{def:b1:lewis-resonance-vsepr:lewis}{par no enlazante}, un S=O y un \ce{S+-O-} en cada una
de las dos \omterm{def:b1:lewis-resonance-vsepr:resonance}{formas resonantes} (o dos S=O con octeto expandido): angular.
\ce{SO3}: trigonal plana, tres enlaces S--O equivalentes. \ce{O3}: O
central con un \omterm{def:b1:lewis-resonance-vsepr:lewis}{par no enlazante}, un O=O y un O--O$^-$, dos \omterm{def:b1:lewis-resonance-vsepr:resonance}{formas
resonantes}: angular.
\end{solution}

\begin{solution}{exo:b1:p-block-16-18:4}
\ce{SO2 + H2O <=> H2SO3}; \ce{SO2 + 2NaOH -> Na2SO3 + H2O};
\ce{SO3 + H2O -> H2SO4}; \ce{SO3 + 2NaOH -> Na2SO4 + H2O}.
\end{solution}

\begin{solution}{exo:b1:p-block-16-18:5}
HF: $x^2/(0.10 - x) = 10^{-3.16}$, $x = \qty{8.0e-3}{mol/L}$, pH 2.10. HCl:
pH 1.00. El enlace H--F es fuerte y el ion fluoruro forma \omterm{def:b1:intermolecular-forces-solvents:hbond}{enlaces de
hidrógeno} fuertes: el HF es un \omterm{def:b1:predominant-reaction:strong-acid}{ácido débil}.
% ledger: dfg:HF_ao, dfg:F-_ao
\end{solution}

\begin{solution}{exo:b1:p-block-16-18:6}
\ce{BrF5}: cinco enlaces y un \omterm{def:b1:lewis-resonance-vsepr:lewis}{par no enlazante}, piramidal cuadrada.
\ce{IF7}: bipiramidal pentagonal. \ce{ICl4-}: cuatro enlaces y dos \omterm{def:b1:lewis-resonance-vsepr:lewis}{pares
no enlazantes}, plano-cuadrada. \ce{XeO3}: tres enlaces y un \omterm{def:b1:lewis-resonance-vsepr:lewis}{par no
enlazante}, piramidal trigonal.
\end{solution}

\begin{solution}{exo:b1:p-block-16-18:7}
$\log K = 2(1.36 - 0.53)/0.059 = 28$. El yodo liberado forma con el
almidón el \omterm{def:b1:complexation:complex}{complejo} azul intenso.
\end{solution}

\begin{solution}{exo:b1:p-block-16-18:8}
$E^\circ(\ce{O3}/\ce{O2}) = 2.07 > 0.53$~V: \ce{O3 + 2I- + 2H+ -> O2 + I2
+ H2O}. El yodo formado se valora con tiosulfato
(\cref{exo:b1:nernst:7}): un \ce{I2} por cada \ce{O3}.
\end{solution}

\begin{solution}{exo:b1:p-block-16-18:9}
\ce{C12H22O11 -> 12C + 11H2O}; $M = \qty{342.0}{g/mol}$, $10.0/342.0 =
\qty{0.0292}{mol}$; carbono, $0.0292 \times 12 \times 12.0 = \qty{4.21}{g}$.
\end{solution}

\begin{solution}{exo:b1:p-block-16-18:10}
$x(0.10 + x)/(0.10 - x) = 10^{-1.99}$: $x = \qty{8.6e-3}{mol/L}$,
$[\ce{H3O+}] = 0.1086$, $\mathrm{pH} = 0.96$ en lugar de 1.00: la segunda
acidez añade alrededor de un \qty{9}{\percent} a $[\ce{H3O+}]$.
\end{solution}

\begin{solution}{exo:b1:p-block-16-18:11}
\ce{F2}/\ce{F-} (2.89~V) está por encima de todos los demás pares: ningún
\omterm{def:b1:nernst:couple}{oxidante} puede arrancarle el electrón al fluoruro. En una electrólisis
acuosa reaccionaría en su lugar el agua, que se oxida a un potencial
mucho más bajo (\ce{O2}/\ce{H2O}, 1.23~V); un fluoruro fundido no
contiene agua.
\end{solution}

\begin{solution}{exo:b1:p-block-16-18:12}
$\Delta_r G^\circ = -51.4 - (16.40 - 51.57) = \qty{-16.2}{kJ/mol}$,
$\log K = 16.2/5.708 = 2.84$, $K = 7 \times 10^2$.
% ledger: dfg:I2_ao, dfg:I-_ao, dfg:I3-_ao
\end{solution}

\begin{solution}{pb:b1:p-block-16-18:1}
\textbf{1.} 0, $+$IV, $+$VI, $+$VI.
\textbf{2.} \ce{S + O2 -> SO2}.
\textbf{3.} S con dos dominios enlazantes y un \omterm{def:b1:lewis-resonance-vsepr:lewis}{par no enlazante}: angular,
de unos
\ang{120}.
\textbf{4.} $10^6/32.1 = \qty{3.12e4}{mol}$.
\textbf{5.} \qty{3.115e4}{mol} de \ce{O2}: $V = 3.115 \times 10^4 \times
8.314 \times 298.15/10^5 = \qty{772}{m^3}$; aire, $772/0.21 = \qty{3.7e3}{m^3}$.
\textbf{6.} La combustión es muy exotérmica; el \omterm{def:b1:elementary-steps:catalyst}{catalizador} funciona en
un intervalo de temperatura limitado, y la oxidación del \ce{SO2},
también exotérmica, es menos completa en caliente.
\textbf{7.} \ce{2SO2 + O2 <=> 2SO3}.
\textbf{8.} Una constante de equilibrio grande no dice nada de la
velocidad: a temperatura ambiente la reacción no avanza.
\textbf{9.} Proporciona un camino de menor \omterm{def:b1:rate-laws:activation-energy}{energía de activación} y se
regenera: un \omterm{def:b1:elementary-steps:catalyst}{catalizador} heterogéneo.
\textbf{10.} Trigonal plana.
\textbf{11.} Más oxígeno mantiene más bajo el cociente
$Q = p_{\ce{SO3}}^2/(p_{\ce{SO2}}^2 p_{\ce{O2}})$: la conversión del
\ce{SO2} avanza más.
\textbf{12.} \qty{0.3}{\percent}.
\textbf{13.} \ce{SO3 + H2SO4 -> H2S2O7}.
\textbf{14.} Forma una niebla de finas gotitas de ácido que atraviesa el
absorbedor.
\textbf{15.} \ce{H2S2O7 + H2O -> 2H2SO4}.
\textbf{16.} \ce{2S + 3O2 + 2H2O -> 2H2SO4}.
\textbf{17.} Una molécula de agua por cada azufre:
$\num{3.115e4} \times 18.0 =
\qty{561}{kg}$.
\textbf{18.} La dilución libera mucho calor; el agua vertida sobre el
ácido hierve de inmediato y salpica ácido.
\textbf{19.} $x(0.050 + x)/(0.050 - x) = 10^{-1.99}$: $x = \qty{7.5e-3}{mol/L}$,
$[\ce{H3O+}] = \qty{0.0575}{mol/L}$, $\mathrm{pH} = 1.24$.
\textbf{20.} $3.12 \times 10^4 \times 0.995 \times 98.1 = \qty{3.04}{t}$.
\textbf{21.} $84 \times 98.1/32.1 = \qty{257}{Mt}$.
\textbf{22.} $98.1/32.1 =$ \textbf{\qty{3.06}{t} de ácido sulfúrico por
tonelada de azufre}.
\end{solution}
